% 第 5 章 Singularities along Curves (原书第 89-100 页)
\chapter{沿曲线的奇性}
\label{chap:5}

\setcounter{mtthm}{-1}

\section{Bessel 函数的渐近性}
\label{sec:5.1}
\index{Bessel 函数!的渐近性 (asymptotics of Bessel functions)}

在本节中, 我们将计算乘子 $\mu_{\lambda}(\xi)=\bigl(1-|\xi|^2\bigr)_+^{\lambda}$ 的 Fourier 变换 $K_{\lambda}$; 我们的目标是估计
\[
K_{\lambda}(x)=K(x)+E(x);
\]
这里
\begin{align*}
K(x)&=C_{\lambda}|x|^{-\frac{n+1+2\lambda}{2}}\cos\Bigl(2\pi|x|-(n+1)\pi-\frac{\lambda\pi}{2}\Bigr), \\
E(x)&=O\Bigl(|x|^{-\frac{n+8+2\lambda}{2}}\Bigr).
\end{align*}
我们遵循 Stein 与 Weiss 的处理方式, 从 Bessel 函数的性质出发详细推导这些估计.

\begin{fact}\label{fact:5.0}
\[
J_{\lambda}(R)=C_{\lambda}R^{\lambda}\int_{-1}^{1}e^{iRt}\bigl(1-t^{2}\bigr)^{\lambda-\frac{1}{2}}\,dt.
\]
\end{fact}

\begin{proof}
这在命题~\ref{prop:2.3} 中已经讨论过.
\end{proof}

\begin{lemma}\label{lem:5.1}
若 $d\theta$ 表示单位球面 $S^{n-1}$ 上的曲面测度, 则
\[
d\check{\theta}(x)=\int_{S^{n-1}}e^{2\pi ix\xi}\,d\theta(\xi)
=c_{n}\frac{J_{\frac{n-2}{2}}\bigl(2\pi|x|\bigr)}{|x|^{\frac{n-2}{2}}}.
\]
\end{lemma}

\begin{proof}
设 $x=|x|x'$, 其中 $|x'|=1$. 我们用低一维的球面来对球面作参数化; 在 $\mathbb{R}^3$ 中这便是平行于赤道而半径各异的圆. 若取北极固定为 $\bar{n}$, 则顶点在单位球面上、且沿这些圆之一的向量 $\bar{\xi}'$ 由条件 $\bar{n}\cdot\bar{\xi}'=\text{常数}$ 刻画. 极点处当然有 $\bar{n}\cdot\bar{\xi}'=\pm 1$; 赤道处 $\bar{n}\cdot\bar{\xi}'=0$.

北极可以任意选取; 为了简单起见, 我们把 $x\in\mathbb{R}^n$ 视为固定的, 并把北极安排为 $\bar{x}'$. 于是对顶点在球面上的 $\bar{\xi}'$, $\bar{x}'$ 与 $\xi'$ 之间的夹角 $\theta$ 属于 $(-\pi,\pi]$. 若令 $L_\theta=\{\bar{\xi}'\in S^{n-1} \mid \bar{x}'\cdot\bar{\xi}'=\cos\theta\}$, 便得到 $S^{n-1}$ 按 $S^{n-2}$ 球面的一个分解. 这是以 $\bar{x}'$ 为北极的球极坐标. 球面 $L_\theta$ 的半径为 $|\sin\theta|$, 故 $L_\theta$ 的体积为 $c_n|\sin\theta|^{n-2}$. 最后我们可以计算:
\begin{align*}
d\check{\theta}(x)&=\int_{-\pi}^{\pi}\int_{L_{\theta}}e^{2\pi i|x|\,\bar{x}'\cdot\bar{\xi}'}\,d\eta\,d\theta \\
&=\int_{-\pi}^{\pi}e^{2\pi i|x|\cos\theta}\int_{L_{\theta}}1\,d\eta\,d\theta
=\int_{-\pi}^{\pi}e^{2\pi i|x|\cos\theta}|L_{\theta}|\,d\theta
=c_{n}\int_{-\pi}^{\pi}e^{2\pi i|x|\cos\theta}\,|\sin\theta|^{n-2}\,d\theta \\
&=c_{n}\int_{-1}^{1}e^{2\pi i|x|y}\bigl(1-y^{2}\bigr)^{\frac{n-3}{2}}\,dy
=c_{n}\frac{J_{\frac{n-2}{2}}(2\pi|x|)}{|x|^{\frac{n-2}{2}}}. \qedhere
\end{align*}
\end{proof}

\begin{proposition}\label{prop:5.2}
\[
\check{\mu}_{\lambda}(x)=K_{\lambda}(x)
=c_{\lambda}\frac{J_{\frac{n}{2}+\lambda}(2\pi|x|)}{|x|^{\frac{n}{2}+\lambda}}.
\]
\end{proposition}

\begin{proof}
对 Bessel 微分方程应用幂级数技巧, 得
\[
J_{\mu}(t)=\sum_{j=0}^{\infty}\frac{(-1)^{j}\bigl(\frac{t}{2}\bigr)^{\mu+j}}{\Gamma(j+\mu+1)},
\]
从而
\begin{align*}
\int_{0}^{1}J_{\mu}(st)\,s^{\mu+1}\bigl(1-s^{2}\bigr)^{\nu}\,ds
&=\sum_{j=0}^{\infty}\frac{(-1)^{j}\bigl(\frac{t}{2}\bigr)^{\mu+j}}{\Gamma(j+\mu+1)}
\int_{0}^{1}s^{\mu+2j}\bigl(1-s^{2}\bigr)^{\nu}s\,ds \\
&=\sum_{j=0}^{\infty}\frac{(-1)^{j}\bigl(\frac{t}{2}\bigr)^{\mu+j}}{\Gamma(j+\mu+1)}\cdot\frac{1}{2}\int_{0}^{1}r^{\mu+j}(1-r)^{\nu}\,dr.
\end{align*}
回忆 Beta 积分, 这就是
\begin{align*}
\frac{1}{2}\sum_{j=0}^{\infty}&\frac{(-1)^{j}\bigl(\frac{t}{2}\bigr)^{\mu+j}}{j!}\,
\frac{\Gamma(\nu+1)\,\Gamma(\mu+j+1)}{\Gamma(\mu+\nu+j+2)\,\Gamma(\mu+j+1)} \\
&=\frac{\Gamma(\nu+1)\,2^{\nu+1}}{t^{\nu+1}}\sum_{j=0}^{\infty}\frac{(-1)^{j}\bigl(\frac{t}{2}\bigr)^{\mu+\nu+1+j}}{\Gamma(j+\nu+\mu+2)}
=c_{\nu}\frac{J_{\mu+\nu+1}(t)}{t^{\nu+1}}.
\end{align*}
在我们的情形, 要计算的是 $K_{\lambda}(x)$, 它是一个 Fourier 变换. 利用极坐标与 $d\check{\theta}(x)$ 的 Bessel 函数表达式, 我们得到
\begin{align*}
K_{\lambda}(x)&=c_{\lambda}\int_{0}^{1}(1-r^{2})^{\lambda}\,r^{n-1}\,
\frac{J_{\frac{n-2}{2}}(2\pi r|x|)}{(r|x|)^{\frac{n-2}{2}}}\,dr \\
&=c_{\lambda}|x|^{\frac{n-2}{2}}(2\pi|x|)^{-(\lambda+1)}\,J_{\frac{n-2}{2}+\lambda+1}(2\pi|x|).
\end{align*}
\end{proof}

\begin{proposition}\label{prop:5.3}
\[
J_{\lambda}(r)=c_{\lambda}r^{-\frac{1}{2}}\cos\Bigl(r-\frac{\pi\lambda}{2}-\frac{\pi}{4}\Bigr)+O\Bigl(r^{-\frac{3}{2}}\Bigr).
\]
\end{proposition}

\begin{proof}
这个课题中的标准技巧, 是利用表达式
\[
r^{-\lambda}J_{\lambda}(r)=c\int_{-1}^{1}e^{irs}\bigl(1-s^{2}\bigr)^{\lambda-\frac{1}{2}}\,ds
\]
中被积函数的解析性来改变围道; 若做得巧妙, 项 $(1-s^2)^{\lambda-\frac{1}{2}}$ 就会变得像 $s^{\lambda-\frac{1}{2}}$ 一样. 这要好得多, 因为我们可以利用齐次性来求出对 $r$ 的依赖关系. 当然, 指数的相位会随围道而改变, 这正是展开式中那些 $\cos$ 项的来历.

那就开始吧. 我们写
\[
\int_{-1}^{1}e^{irs}\bigl(1-s^{2}\bigr)^{\lambda-\frac{1}{2}}\,ds
=\lim_{\epsilon\to0}\int_{-1+\epsilon}^{1-\epsilon}e^{irs}\bigl(1-s^{2}\bigr)^{\lambda-\frac{1}{2}}\,ds,
\]
这使我们避开了非解析点. 于是 $e^{irz}(1-z^2)^{\lambda-\frac{1}{2}}$ 在单连通区域 $\mathbb{C}-\{(-\infty,-1+\epsilon]\cup[1-\epsilon,\infty)\}$ 上有一个解析分支. 下面我们会稍微马虎地对待这些 $\epsilon$.

我们取一个高为 $a$、底为 $[-1,1]$ 的矩形, 并沿逆时针方向走过其边界. 边界曲线 $\gamma$ 有四段:
\begin{gather*}
\gamma_{1}=\{s+i0 \mid -1\leq s\leq 1\};\qquad \gamma_{2}=\{1+iy \mid 0\leq y\leq a\}; \\
\gamma_{3}=\{s+ia \mid -1\leq s\leq 1\};\qquad \gamma_{4}=\{-1+i(a-y) \mid 0\leq y\leq a\}.
\end{gather*}
则
\[
\int_{\gamma_{1}}=-\int_{\gamma_{2}}-\int_{\gamma_{3}}-\int_{\gamma_{4}},
\]
我们取 $a\to\infty$ 的极限. 此时
\[
\int_{\gamma_{3}}=-\int_{-1}^{1}e^{-ar}e^{irs}\bigl(1-(s+ia)^{2}\bigr)^{\lambda-\frac{1}{2}}\,ds,
\]
指数因子保证了它收敛于零. 对 $\gamma_{2}$,
\begin{align*}
\bigl(1-(1+iy)^{2}\bigr)^{\lambda-\frac{1}{2}}
&=\bigl(y^{2}-2iy\bigr)^{\lambda-\frac{1}{2}}
=y^{\lambda-\frac{1}{2}}(y-2i)^{\lambda-\frac{1}{2}},
\end{align*}
对 $\gamma_{4}$ 类似. 这两个积分合起来给出
\[
\int_{0}^{\infty}e^{ir(1+y)}\,y^{\lambda-\frac{1}{2}}(y-2i)^{\lambda-\frac{1}{2}}\,dy
-\int_{0}^{\infty}e^{ir(-1+y)}\,y^{\lambda-\frac{1}{2}}(y+2i)^{\lambda-\frac{1}{2}}\,dy.
\]
现在我们把 $(y+2i)^{\lambda-\frac{1}{2}}$ 展开为幂级数, 其首项为 $(2i)^{\lambda-\frac{1}{2}}$; 在每个被积函数中它贡献一项形如 $2^{\lambda-\frac{1}{2}}e^{(\lambda-\frac{1}{2})i\pi/2}$ 的量; 两个积分合起来产生
\[
\cos\Bigl(r-\frac{\pi\lambda}{2}-\frac{\pi}{4}\Bigr)\int_{0}^{\infty}e^{-ry}\,y^{\frac{1}{2}+\lambda}\,dy
=\Gamma\bigl(\lambda+\tfrac{1}{2}\bigr)\,r^{-(\frac{1}{2}+\lambda)}\cos\Bigl(r-\frac{\pi\lambda}{2}-\frac{\pi}{4}\Bigr).
\]
如果我们能控制由误差 $(y+2i)^{\lambda-\frac{1}{2}}-(2i)^{\lambda-\frac{1}{2}}$ 产生的项, 证明即告完成. 若 $y\leq 1$, 我们利用光滑性, 均值定理给出这个差被 $Cy$ 控制. 若 $y>1$, 我们提出因子 $y$ 并再次使用均值定理. 于是
\begin{align*}
\int_{0}^{\infty}&E(y)\,e^{-ry}\,y^{\lambda-\frac{1}{2}}\,dy
\leq\int_{0}^{1}e^{-ry}y^{\lambda+\frac{1}{2}}\,dy+\int_{1}^{\infty}e^{-ry}y^{2\lambda-1}\,dy \\
&=r^{-(\lambda+\frac{3}{2})}\int_{0}^{r}e^{-y}y^{\lambda+\frac{1}{2}}\,dy
+r^{-2\lambda}\int_{r}^{\infty}e^{-y}y^{2\lambda+1}\,dy \\
&=r^{-(\lambda+\frac{3}{2})}\,\Gamma\bigl(\lambda+\tfrac{3}{2}\bigr)
+C_{k}\,r^{-2\lambda}\int_{r}^{\infty}y^{-k}y^{2\lambda+1}\,dy
=O\bigl(r^{-(\lambda+\frac{3}{2})}\bigr). \qedhere
\end{align*}
\end{proof}

\begin{corollary}\label{cor:5.4}
当 $\lambda>\frac{n-1}{2}$ 时, $T_\lambda$ 在 $L^p(\mathbb{R}^n)$ ($1\leq p\leq\infty$) 上有界; 而当 $\lambda<\frac{n-1}{2}$ 且 $p$ 不在范围
\[
\frac{2n}{n+1+2\lambda}<p<\frac{2n}{n-1-2\lambda}
\]
之内时, $T_\lambda$ 在 $L^p(\mathbb{R}^n)$ 上无界 (对照图~\ref{fig:2}).
\end{corollary}

\begin{proof}
若 $\lambda>\frac{n-1}{2}$, 则 $K_\lambda\in L^1$. 为得到无界性结果, 我们想取 $f$ 为一个 $\delta$ 函数; 渐近性将会表明
\[
K_{\lambda}\star f\approx|x|^{-\frac{n+1+2\lambda}{2}},
\]
而除非 $\frac{2n}{n+1+2\lambda}<p$, 这不可能属于 $L^p$. 由于 $\delta$ 函数不在 $L^p$ 中, 我们必须更巧妙一些; 关键在于取一个支集非常小的 $f$, 使得卷积核中的振荡可以被抑制. 若把 $f$ 想成紧支集的, 卷积就是把 $f$ 对 $K_\lambda$ 的平移作平均, 于是有时这个平均几乎为零. 我们必须分析 $f$ 与 $K_\lambda$ 的峰相遇的那些位置.

设 $D$ 是半径为 $\frac{1}{100}$ 的球, 并设 $A_{k}$ 表示环形区域
\[
\{x \mid \bigl|2\pi|x|-\lambda\pi-\frac{\pi}{2}-2k\pi\bigr|\leq 10^{-2}\}.
\]
注意 $A_{k}$ 的测度约为 $Ck^{n-1}$, 且若 $x\in A_{k}$, $y\in D$, 则
\[
\cos\Bigl(2\pi|x-y|-(n+1)\pi-\frac{1}{2}\bigl(\frac{n}{2}+\lambda\bigr)\pi\Bigr)\geq\frac{1}{2}.
\]
于是, 若取 $f=\chi_{D}$,
\[
\|T_{\lambda}f\|_{p}\geq\bigl\|(K_{\lambda}\star f)\,\chi_{\bigcup A_{k}}\bigr\|_{p}
-\bigl\|(E\star f)\,\chi_{\bigcup A_{k}}\bigr\|_{p}.
\]
但
\begin{align*}
(K_{\lambda}\star f)(x)\,\chi_{\bigcup A_{k}}(x)
&=\int_{D}\frac{\cos\bigl(2\pi|x-y|+\mathrm{mess}\bigr)}{|x-y|^{\frac{n+1+2\lambda}{2}}}\,dy
\geq C\,k^{-(\frac{n+1+2\lambda}{2})};
\end{align*}
类似地,
\[
|E\star f|\,\chi_{\bigcup A_{k}}\leq C\,k^{-(\frac{n+3+2\lambda}{2})}.
\]
最后,
\[
\|T_{\lambda}f\|_{p}\geq C\left(\sum_{k}k^{-\frac{n+1+2\lambda}{2}}|A_{k}|\right)^{\frac{1}{p}}.
\]
这个级数有限当且仅当 $(n-1)-(n+1+2\lambda)\frac{p}{2}<-1$.
\end{proof}

\begin{remark}\label{rem:5.5}
我们已经看到, 沿曲线的奇性比沿直线的奇性糟糕得多. 单位正方形的特征函数的卷积核是
\[
\frac{\sin x_{1}}{x_{1}}\cdot\frac{\sin x_{2}}{x_{2}},
\]
它对所有 $p>1$ 属于 $L^{p}$. 与之相对, 单位圆盘的特征函数的变换是
\[
\frac{J_{1}(|x|)}{|x|},
\]
它仅在 $p>\frac{4}{3}$ 时属于 $L^p$. 这两个结果的详细证明我们都已看到; 但隐藏在这一区别之下的几何却仍让我们困惑.

进入 $L^p$ 的标准途径是从 $L^1$ 与 $L^2$ 内插; 由于对小的 $\lambda$, $T_\lambda$ 在 $L^1$ 上无界, 这条路走不通. 我们求助于标准模型: 把光滑紧支集函数 $\psi_j$ 扩张到 $\mathbb{R}^n$, 办法是从定义在 $\mathbb{R}^1$ 上的 $\psi_0$ 出发, 令 $\psi_j(\xi)=\psi_0\bigl(2^j(|\xi|-1)\bigr)$. 这些 $\psi_j$ 的支集在一个厚度约 $2^{-j}$、以 $|\xi|=1$ 为中心的环形区域中. (这与一维的图景颇为不同, 但请记住, 我们希望解决的奇性正在 $|\xi|=1$ 处.) 我们预计 $\mu_\lambda\approx 2^{-j\lambda}$, 且 ${}_p\|T_\lambda\|\leq\sum 2^{-j\lambda}\,{}_p\|T_j\|$. 然而与一维情形不同, 这些 $\psi_j$ 并不是同一个函数的伸缩, 因此诸 $T_j$ 的算子范数并不相等. 为了控制这些范数, 我们考察 $L^1$ 与 $L^2$ 上的范数并作内插.

在 $L^2$ 上, ${}_2\|T_j\|=\|\psi_j\|_\infty=1$. 在 $L^1$ 上, ${}_1\|T_j\|=\|\check{\psi}_j\|_1$. 直观上, $\check{\psi}_j$ 恰是圆盘特征函数的 Fourier 变换限制在环形区域 $\{2^j<|x|<2^{j+1}\}$ 上的部分. 一个快速的计算表明我们预计 $\|\check{\psi}_j\|_1=C2^{\frac{n-1}{2}j}$. 利用渐近性可以把这个算得更精确:
\[
\check{\psi}_{j}(x)=\int\frac{J_{\frac{n-2}{2}}(2\pi|x|r)}{(|x|r)^{\frac{n-2}{2}}}\,\psi\bigl(2^{j}(|r|-1)\bigr)\,r^{n-1}\,dr.
\]
在区域 $|x|<1$ 中, $|\check{\psi}_j|\leq\|\psi_j\|_1\leq 1$; 在区域 $|x|>1$ 中, 对那些 $r\in\operatorname{support}\psi_j$, $r$ 接近 $1$, 于是 $r|x|>1$, 即我们处于 Bessel 函数渐近性成立的区域. 主项贡献了
\[
C|x|^{-\frac{n-2}{2}}\int e^{2\pi ir|x|}\,\psi\bigl(2^{j}(|r|-1)\bigr)\,r^{\frac{n-1}{2}}\,dr.
\]
令 $s=2^{j}(|r|-1)$; 我们得到
\[
C|x|^{-\frac{n-2}{2}}e^{2\pi i|x|}2^{-j}\int e^{2\pi is2^{-j}|x|}\,\psi(s)\,(1+s2^{-j})^{\frac{n-1}{2}}\,ds.
\]
若把 $(1+s2^{-j})^{\frac{n-1}{2}}$ 在 $s=0$ 附近展开, 主项是
\begin{align*}
C|x|^{-\frac{n-2}{2}}e^{2\pi i|x|}2^{-j}\int e^{2\pi is2^{-j}|x|}\,\psi(s)\,ds
&=C|x|^{-\frac{n-2}{2}}e^{2\pi i|x|}2^{-j}\,\check{\psi}_{0}(2^{-j}|x|).
\end{align*}
注意到 $\psi_0\in\mathcal{S}$; 总之, 若令
\[
\varphi(x)=|x|^{-\frac{n-1}{2}}\check{\psi}_{0}(|x|),
\]
则 $\varphi\in L^1$, 且
\[
|\check{\psi}_{j}(x)|\leq C\,2^{\frac{n-1}{2}j}\,2^{-jn}\,\varphi(2^{-j}x).
\]
于是 $\|\check{\psi}_j\|_1\leq C2^{\frac{n-1}{2}j}\|\varphi\|_1=C2^{\frac{n-1}{2}j}$. 这印证了我们的直观, 但应当提醒读者: 在上面的计算中, 我们对误差项的控制是非常马虎的.

现在可以内插了: 若 $\frac{1}{p}=\frac{t}{1}+\frac{1-t}{2}$, 即 $t=\frac{2}{p}-1$, 则
\begin{align*}
{}_p\|T_{j}\|&\leq{}_{1}\|T_{j}\|^{t}\,{}_{2}\|T_{j}\|^{1-t}
\leq 2^{j\left(\frac{n-1}{2}\right)t}\cdot 1^{t}\leq 2^{j(n-1)\left(\frac{1}{p}-\frac{1}{2}\right)}.
\end{align*}
于是, 当 $p\leq 2$ 时, 只要
\[
\lambda>(n-1)\left(\frac{1}{p}-\frac{1}{2}\right),
\]
我们就能使几何级数求和.
\end{remark}

\begin{figure}[htbp]
\centering
\begin{tikzpicture}[scale=0.9]
  % 阴影区域 (竖线填充)
  \fill[pattern=vertical lines, pattern color=black!70]
    (0,3.4) -- (2.5,0) -- (5,3.4) -- (5,4.2) -- (0,4.2) -- cycle;
  % 边框
  \draw (0,0) rectangle (5,4.2);
  % 虚线 V 形边界
  \draw[dashed] (0,3.4) -- (2.5,0);
  \draw[dashed] (2.5,0) -- (5,3.4);
  % 标记点
  \draw (0,3.4) circle (2.4pt);
  \draw (5,3.4) circle (2.4pt);
  \fill (2.5,0) circle (2.4pt);
  % 标注
  \node[left] at (-0.1,3.4) {$\alpha=\frac{n-1}{2}$};
\end{tikzpicture}
\caption{Bessel 函数的振荡与截断}
\label{fig:3}
\end{figure}

\section{de Leeuw 与 Stein 的一个定理}
\label{sec:5.2}

\begin{theorem}\label{thm:5.6}
若 $0<\lambda<\frac{n-1}{2}$, 则当 $\lambda>(n-1)\bigl|\frac{1}{p}-\frac{1}{2}\bigr|$ 时, $T_\lambda$ 在 $L^p(\mathbb{R}^n)$ 上有界 (对照图~\ref{fig:3}).
\end{theorem}

\begin{proof}
当然, 我们可以沿标准模型的轮廓给出一个证明; 不过为了展示技术的多样性, 我们将给出一个对 $K_{\lambda}$ (而非 $\mu_{\lambda}$) 作分解的证明.

标准分解提示我们把 $K_\lambda$ 分解为支在二进环形区域上的小块. 为此, 设 $K\in\mathcal{S}$ 满足 $0\leq K\leq 1$, 当 $|x|\leq 100$ 时 $K(x)=1$, 当 $|x|>200$ 时 $K(x)=0$. 令 $\phi_j(x)=K\bigl(\frac{|x|}{2^j}\bigr)-K\bigl(\frac{|x|}{2^{j+1}}\bigr)$. 注意 $\phi_j$ 支在环形区域 $100\cdot 2^j<|x|<200\cdot 2^j$ 上, 于是 $\phi_j$ 的作用如同某个环形区域特征函数的光滑化版本.

现在令 $K_j=\phi_jK_\lambda$. 则 $\sum_j\phi_j(x)=K(x)-K(0)$, 于是 $K(x)-\sum_j\phi_j(x)=K(0)=1$; 因此 $K_\lambda K(x)=\sum_jK_j+K_\lambda$, 从而 $K_\lambda$ 与 $\sum_jK_j$ 只相差 $K_\lambda\bigl(1-K(x)\bigr)$. 由于 $K_\lambda\bigl(1-K(x)\bigr)$ 有界且紧支集, 它给出一个 $L^1$ 误差项, 在所有 $L^p$ 上是有界算子.

在这幅图景中, $\|K_j\|_1$ 容易计算, 而 $\|\hat{K}_j\|_\infty$ 却难. 前者给出:
\[
\int\left|K_{\lambda}(x)\right|\,\left|\phi_{j}(x)\right|dx
\leq C\int_{100\cdot 2^{j}}^{200\cdot 2^{j}}r^{-\frac{n+1+2\lambda}{2}}\,r^{n-1}\,dr
=C\,2^{-\lambda j}\,2^{j\frac{n-1}{2}},
\]
与之前一样. 为了对 $\|\hat{K}_j\|_\infty$ 可能有多大有一个直观感受, 注意 $\hat{K}_j=\hat{K}_\lambda\star\hat{\phi}_j$, 而
\[
\hat{\phi}_{j}(\xi)=2^{nj}\hat{K}(2^{j}\xi)-2^{n(j-1)}\hat{K}(2^{j-1}\xi).
\]
若令 $\delta=2^{-j}$, $\psi(\xi)=2\hat{K}(\xi)-\hat{K}\bigl(\frac{\xi}{2}\bigr)$, 则可以把这些改写为
\[
\psi_{\delta}(\xi)=\delta^{-n}\psi(\delta^{-1}\xi)=\hat{\phi}_{j}(\xi).
\]
现在 $\psi_\delta$ 是一个标准的恒等逼近, 只是 $\int\psi_\delta=0$. 当 $\delta$ 变小, 即 $j\to\infty$ 时, 我们预计 $\psi_\delta\star\mu_\lambda$ 将收敛到 $\bigl(\int\psi_\delta\bigr)\mu_\lambda=0$, 并且预计收敛速度由 $\mu_\lambda$ 的光滑性控制. 于是我们预计有估计:
\[
|\psi_{\delta}\star\mu_{\lambda}-0|\leq C\delta^{\lambda}=C2^{-j\lambda}.
\]
有一种更直观的思考方式: 把 $\psi$ 想成一个圆盘的特征函数; 那么 $\psi_{\delta}\star f$ 就是在一个更小的圆盘上对 $f$ 作平均. 在 $|\xi|=1$ 附近, $\mu_{\lambda}$ 的最大值是 $C\delta^{\lambda}$, 它支配着这个平均. 远离 $|\xi|=1$ 处, $\mu_{\lambda}$ 光滑, 于是这个平均如同 $\mu_{\lambda}\int\psi_{\delta}=0$.

直观分析至此完成; 如同本节开头的论证, 估计
\[
\|K_{j}\|_{1}\leq C\,2^{j\frac{n-1}{2}}\,2^{-j\lambda};\qquad
\|\hat{K}_{j}\|_{\infty}\leq C2^{-j\lambda}
\]
足以通过内插完成定理的证明. 我们现在的义务, 是给这些估计背后的直观一个严格 (尽管乏味) 的解读. 不过这个证明会引出几个重要的技巧, 所以它是一定要看的.

在区域 $|\xi|>2$ 中, $\psi_\delta$ 与 $\mu_\lambda$ 的``支集''不相交, 于是 $\psi_\delta\star\mu_\lambda$ 应当为零. 由于 $\psi_\delta$ 实际上并不紧支集, 我们需要另一个窍门. $\psi\in\mathcal{S}$ 的急速下降应当能解决问题. 于是,
\[
|\psi_{\delta}(\eta)|=|\delta^{-n}\psi(\delta^{-1}\eta)|
\leq C\delta^{-n}\,|\delta^{-1}\eta|^{-(n+i)}
=C\delta^{i}\,|\eta|^{-(n+i)},
\]
\[
|K_{j}(\xi)|\leq\int|\mu_{\lambda}(\xi-\eta)|\,|\psi_{\delta}(\eta)|\,d\eta
\leq C\|\mu_{\lambda}\|_{\infty}\,\delta^{i}\int\chi_{D}(\xi-\eta)\,|\eta|^{-(n+i)}\,d\eta;
\]
这里 $D=\operatorname{support}\mu_\lambda$. 由于必须有 $|\xi-\eta|\leq 1$, 有 $|\xi|-|\eta|<1$, 即 $|\eta|>|\xi|-1>1$, 于是这个积分被 $C$ 控制, 取 $i$ 充分大便可使其任意小.

在区域 $|\xi|<2$ 中, 我们预计可以用 $\psi_\delta\star\mu_\lambda$ 收敛到零的速度来控制其大小. 典型地, 我们会估计
\begin{align*}
\left|\psi_{\delta}\star\mu_{\lambda}(\xi)-\Bigl(\int\psi_{\delta}\Bigr)\mu_{\lambda}(\xi)\right|
&\leq\int|\mu_{\lambda}(\xi-\delta\eta)-\mu_{\lambda}(\xi)|\,|\psi(\eta)|\,d\eta.
\end{align*}
现在我们想用估计
\[
|\mu_{\lambda}(\xi-\delta\eta)-\mu_{\lambda}(\xi)|\leq C|\delta\eta|^{\lambda}
\]
来完成证明. 不幸的是, 当 $\lambda>1$ 时这个估计是错的; 不过, 若 $\lambda\leq 1$, 我们可以完成证明:
\begin{align*}
\int|\mu_{\lambda}(\xi-\delta\eta)-\mu_{\lambda}(\xi)|\,|\psi(\eta)|\,d\eta
&\leq C\int|\delta\eta|^{\lambda}\,|\psi(\eta)|\,d\eta\leq C\delta^{\lambda}.
\end{align*}
若 $\lambda>1$, 我们需要分部积分 $k=[\lambda]$ 次; 这将把 $\mu_\lambda$ 微分成一团乱麻乘以 $(1-|\xi|^2)^{\lambda-k}$, 而此时 $\lambda-k\leq 1$. 为使奏效, $\psi$ 的 $k$ 次积分 (称之为 $\phi$) 必须具有与 $\psi$ 相同的结构, 即 $\phi\in\mathcal{S}$ 且 $\int\phi=0$. 当然, 对任意的 $\psi$, 它的不定积分一般不满足这一点; 我们必须小心. 如下定义 $\phi$:
\[
\check{\phi}(x)=\frac{\check{\psi}(x)}{(2\pi ix_{1})\cdots(2\pi ix_{k})}.
\]
现在 $\tilde{\psi}=K_1$ 在原点的一个邻域中消失, 于是 $\tilde{\phi}\in\mathcal{S}$, $\phi\in\mathcal{S}$; 且 $\int\phi=\check{\phi}(0)=0$. 此外, $D_{\xi_1}\cdots D_{\xi_k}\phi=\psi$, 并且
\[
\int\eta_{1}\cdots\eta_{k}\,\phi(\eta)\,d\eta
=\bigl(D_{x_{1}}\cdots D_{x_{k}}\check{\phi}\bigr)(0)=0.
\]
再者,
\[
D_{\eta_{1}}\cdots D_{\eta_{k}}\,\mu_{\lambda}(\xi-\delta\eta)
=C\,(\delta\eta_{1})\cdots(\delta\eta_{k})\,\mu_{\lambda-k}(\xi-\delta\eta).
\]
于是我们可以分部积分 $k$ 次而不产生任何边界项, 因为 $\mu_{\lambda-k}$ 与 $\phi$ 都在无穷远处消失. 我们要指出, 若诸 $\eta$ 并非互不相同, 这些公式并不准确; 另一方面, 我们只需处理 $\lambda<\frac{n-1}{2}<n$ 的情形, 所以我们有用之不竭的余地来选择互不相同的 $\eta$. 为完成证明:
\begin{align*}
\bigl|\hat{K}_{j}(\xi)\bigr|&=\left|\int\mu_{\lambda}(\xi-\delta\eta)\,\psi(\eta)\,d\eta\right| \\
&=C\delta^{k}\left|\int\mu_{\lambda-k}(\xi-\delta\eta)\,\eta_{1}\cdots\eta_{k}\,\phi(\eta)\,d\eta
-\mu_{\lambda-k}\int\eta_{1}\cdots\eta_{k}\,\phi(\eta)\,d\eta\right| \\
&\leq C\delta^{k}\int|\mu_{\lambda-k}(\eta-\delta\eta)-\mu_{\lambda-k}(\eta)|\,\bigl|\eta_{1}\cdots\eta_{k}\bigr|\,|\phi(\eta)|\,d\eta \\
&\leq C\delta^{k}\,\delta^{\lambda-k}\int|\eta|^{\lambda-k}\,|\eta|^{k}\,|\phi(\eta)|\,d\eta
\leq C\delta^{\lambda}. \qedhere
\end{align*}
\end{proof}

看一眼图~\ref{fig:3} 中的直线段便可知发生过某种内插. 但那是哪种内插呢? 毕竟, 图~\ref{fig:3} 画的并不是某个算子的 $L^p$--$L^q$ 有界性, 而是一族算子的 $L^p$--$L^p$ 有界性. 要作内插, 需要 ${}_p\|T_\lambda\|$ 作为 $\lambda$ 的函数具有某种凸性. 如果研究一下 Riesz--Thorin 定理的证明, 其要求是 $F(z)=\int T_zfg$ 在带形区域上解析、在边界上连续、且增长不快于 $Ce^{-b|y|}$. 我们把这样的算子族称为\emph{容许} (admissible) 的.

\begin{theorem}\label{thm:5.7}
\index{内插!复内插 (complex interpolation)}
\index{容许算子族 (admissible family)}
设 $T_{z}$ 是区域 $0\leq\operatorname{Re}z\leq 1$ 中的一个容许算子族. 设 $T_{0+iy}$ 在 $L^{p_0}$ 上有界, $T_{1+iy}$ 在 $L^{p_1}$ 上有界. 若
\[
{}_{p_{j}}\|T_{j+iy}\|\leq e^{C\,e^{-b|y|}}
\]
对某个 $b<\pi$ 成立, 则对 $\frac{1}{p}=\frac{t}{p_0}+\frac{1-t}{p_1}$, $T_t$ 在 $L^p$ 上有界.
\end{theorem}

\begin{proof}
这个称为复内插的结果归功于 Stein. 尽管这是一个重要的结果, 详细的证明会使我们离题太远; 参见章末的文献.

为证明该定理, Stein 首先把三线定理推广, 得到估计
\[
\log|F(z)|\leq\frac{1}{2}\sin\pi x\int_{-\infty}^{\infty}
\left[\frac{\log|F(iy)|}{\cosh\pi y-\cos\pi x}
+\frac{\log|F(1+iy)|}{\cosh\pi y+\cos\pi x}\right]dy.
\]
Stein 现在取 $f=\sum a_j\chi_{E_j}$, $g=\sum b_k\chi_{F_k}$ 为满足 $\|f\|_p=\|g\|_{p'}=1$ 的简单函数. 使用定理~\ref{thm:1.13} 的记号, 令
\[
F(z)=\int T_{z}\left(\sum|a_{j}|^{p\alpha}e^{i\theta_{j}}\chi_{E_{j}}\right)\left(\sum|b_{k}|^{p'(1-\alpha)}e^{i\psi_{k}}\chi_{F_{k}}\right),
\]
其中 $\alpha(z)=\frac{1-z}{p_0}+\frac{z}{p_1}$. 则与定理~\ref{thm:1.13} 的证明一样,
\begin{align*}
|F(1+iy)|&\leq{}_{p_{1}}\|T_{1+iy}\|\leq M_{1}(y), \\
|F(0+iy)|&\leq{}_{p_{0}}\|T_{0+iy}\|\leq M_{0}(y),
\end{align*}
这里诸 $M_j(y)$ 的增长不快于 $e^{Ce^{-b|y|}}$. 于是修改后的三线定理给出
\begin{align*}
|F(z)|\leq\exp\frac{1}{2}\sin\pi x\int_{-\infty}^{\infty}
&\frac{\log\bigl|M_{0}(y)\bigr|}{\cosh\pi y-\cos\pi x} \\
&+\frac{\log\bigl|M_{1}(y)\bigr|}{\cosh\pi y+\cos\pi x}\,dy,
\end{align*}
而 $M_j(y)$ 的增长保证了积分的收敛性. 于是 $|F(z)|\leq C_t$, 正如所需.
\end{proof}

\begin{corollary}\label{cor:5.8}
设 $0<\lambda<\frac{n-1}{2}$. 则当 $\lambda>(n-1)\bigl|\frac{1}{p}-\frac{1}{2}\bigr|$ 时, $T_\lambda$ 在 $L^p$ 上有界.
\end{corollary}

\begin{proof}
令
\[
K_{z}(x)=c_{z}\frac{J_{\frac{n}{2}+z}(2\pi|x|)}{|x|^{\frac{n}{2}+z}},
\]
其中 $c_{z}$ 待定. 为检验容许性,
\[
\int T_{z}(f)\,g\leq\|T_{z}f\|_{2}\,\|g\|_{2}\leq\|f\|_{2}\,\|g\|_{2}.
\]
类似地,
\[
{}_{2}\|T_{0+iy}\|=\|\mu_{0+iy}\|_{\infty}\leq 1.
\]
但
\[
{}_{1}\|T_{1+iy}\|\leq\|K_{\lambda+iy}\|_{1}.
\]
现在归一化常数变得重要了! 由引理~\ref{lem:5.1} 与命题~\ref{prop:5.2}, Bessel 函数的标准归一化涉及 Gamma 函数的商. 若考察 Gamma 函数作为积分的定义, 就会看到它在复平面上呈指数增长. Bessel 函数本身也有类似的增长; 只需看看它的幂级数展开. 因此平衡这两种增长十分重要; 实际的工作见 Stein~\cite{stein-interpolation}.
\end{proof}

\section{评注与推广}
\label{sec:5.3}

\noindent (5.4):\quad 这里的简单反例之所以有效, 只是因为核的振荡
\[
\cos(2\pi|x|)
\]
是缓慢的; 它对表面相似而振荡为 $\cos(2\pi|x|^2)$ 的核就不会奏效. 这类核的 $L^p$ 行为非常不同; 关于 $\frac{e^{i|x|^\alpha}}{|x|^\beta}$ 的处理, 见 C.~Fefferman~\cite{fefferman-stronglysingular} 与 S.~Wainger~\cite{wainger-special}.

\noindent (5.5):\quad 我们将在第~\ref{chap:6} 章第~1 节中重新着手 $K_{\lambda}$ 的局部计算.

\noindent (5.6):\quad 这里给出的证明归功于 de Leeuw~\cite{deleeuw}.

\noindent (5.7):\quad 这个结果归功于 Stein~\cite{stein-interpolation}; 关于其他内插技术, 见 Stein 与 Weiss~\cite{steinweiss} 第~209 页的讨论.

\noindent (5.8):\quad 这个结果归功于 Stein~\cite{stein-interpolation}.
